\documentclass[12pt, a4paper]{article}

\usepackage[utf8]{inputenc}
\usepackage[T1]{fontenc}
\usepackage[mathscr]{euscript}
\let\euscr\mathscr \let\mathscr\relax
\usepackage[scr]{rsfso}
\usepackage{amssymb,amsmath,amsthm,amsfonts}
\usepackage[shortlabels]{enumitem}
\usepackage{multicol,multirow}
\usepackage{lipsum}
\usepackage{balance}
\usepackage{calc}
\usepackage[colorlinks=true,citecolor=blue,linkcolor=blue]{hyperref}
\usepackage{import}
\usepackage{xifthen}
\usepackage{pdfpages}
\usepackage{transparent}
\usepackage{tabularx}
\usepackage{listings}
\usepackage[dvipsnames]{xcolor}
\usepackage{bbm}

\newcommand{\incfig}[2][1.0]{
	\def\svgwidth{#1\columnwidth}
	\import{./figures/}{#2.pdf_tex}
}
\newcommand{\incimg}[2][1.0]{
	\includegraphics[width=#1\columnwidth]{./figures/#2}
}

\newcommand{\xor}{\oplus}
\newcommand{\floor}[1]{\left\lfloor #1 \right\rfloor}
\newcommand{\ceil}[1]{\left\lceil #1 \right\rceil}

\newlist{enumproof}{enumerate}{4}
\setlist[enumproof,1]{label=\arabic*., parsep=1em}
\setlist[enumproof,2]{label=\arabic{enumproofi}.\arabic*., parsep=1em}
\setlist[enumproof,3]{label=\arabic{enumproofi}.\arabic{enumproofii}.\arabic*., parsep=1em}
\setlist[enumproof,4]{label=\arabic{enumproofi}.\arabic{enumproofii}.\arabic{enumproofiii}.\arabic*., parsep=1em}

\lstdefinestyle{mystyle}{
	language=Python, % Set the language to Python
	commentstyle=\color{codegray}\itshape, % Python comments (#)
	keywordstyle=\color{blue},  % Python reserved keywords
	stringstyle=\color{mauve},             % Strings
	basicstyle={\ttfamily\footnotesize},   % Monospace font
	breakatwhitespace=true,
	breaklines=true,
	captionpos=b,
	keepspaces=true,
	numbers=none,
	tabsize=2,
	showstringspaces=false,
	% Extra keyword groups for Python
	morekeywords=[2]{self, cls, __init__, __str__, __repr__},
	keywordstyle=[2]\color{identblue},     % Special identifiers / methods
	morekeywords=[3]{int, float, bool, str, list, dict, set, tuple, object, None, True, False},
	keywordstyle=[3]\color{codegreen},     % Built-in types
	morekeywords=[4]{print, input, open, range, len, type, isinstance, super, enumerate, zip, map, filter, sum, min, max, sorted},
	keywordstyle=[4]\color{cyan},          % Built-in functions
}
\lstset{style=mystyle}

\usepackage{ifthen}
\usepackage[landscape]{geometry}
\usepackage[shortlabels]{enumitem}

\ifthenelse{\lengthtest { \paperwidth = 11in}}
    { \geometry{top=.5in,left=.5in,right=.5in,bottom=.5in} }
	{\ifthenelse{ \lengthtest{ \paperwidth = 297mm}}
		{\geometry{top=1cm,left=1cm,right=1cm,bottom=1cm} }
		{\geometry{top=1cm,left=1cm,right=1cm,bottom=1cm} }
	}

\pagestyle{empty}
\makeatletter
\renewcommand\thesection{\arabic{section}.}
\renewcommand{\section}{\@startsection{section}{1}{0mm}%
                                {-1ex plus -.5ex minus -.2ex}%
                                {0.05ex}%x
                                {\normalfont\normalsize\bfseries}}
\renewcommand{\subsection}{\@startsection{subsection}{2}{0mm}%
                                {-1ex plus -.5ex minus -.2ex}%
                                {0.05ex}%
                                {\normalfont\small\bfseries}}
\renewcommand{\subsubsection}{\@startsection{subsubsection}{3}{0mm}%
                                {-1ex plus -.5ex minus -.2ex}%
                                {0.05ex}%
                                {\normalfont\footnotesize\bfseries}}
\newcommand{\colbreak}{\vfill\null\columnbreak}
\makeatother
\setcounter{secnumdepth}{1}
\setlength{\parindent}{0pt}
\setlength{\parskip}{0.7em}

\setlist[itemize]{itemsep=0.6ex, topsep=-2pt, partopsep=0pt, parsep=0pt}
\setlist[enumerate]{itemsep=0.6ex, topsep=-2pt, partopsep=0pt, parsep=0pt}

% Things Lie
\newcommand{\kb}{\mathfrak b}
\newcommand{\kg}{\mathfrak g}
\newcommand{\kh}{\mathfrak h}
\newcommand{\kn}{\mathfrak n}
\newcommand{\ku}{\mathfrak u}
\newcommand{\kz}{\mathfrak z}
\DeclareMathOperator{\Ext}{Ext} % Ext functor
\DeclareMathOperator{\Tor}{Tor} % Tor functor
\newcommand{\gl}{\opname{\mathfrak{gl}}} % frak gl group
\renewcommand{\sl}{\opname{\mathfrak{sl}}} % frak sl group chktex 6

% More script letters etc.
\newcommand{\SA}{\mathcal A}
\newcommand{\SB}{\mathcal B}
\newcommand{\SC}{\mathcal C}
\newcommand{\SF}{\mathcal F}
\newcommand{\SG}{\mathcal G}
\newcommand{\SH}{\mathcal H}
\newcommand{\OO}{\mathcal O}

\newcommand{\SCA}{\mathscr A}
\newcommand{\SCB}{\mathscr B}
\newcommand{\SCC}{\mathscr C}
\newcommand{\SCD}{\mathscr D}
\newcommand{\SCE}{\mathscr E}
\newcommand{\SCF}{\mathscr F}
\newcommand{\SCG}{\mathscr G}
\newcommand{\SCH}{\mathscr H}

% Mathfrak primes
\newcommand{\km}{\mathfrak m}
\newcommand{\kp}{\mathfrak p}
\newcommand{\kq}{\mathfrak q}

% number sets
\newcommand{\RR}[1][]{\ensuremath{\ifstrempty{#1}{\mathbb{R}}{\mathbb{R}^{#1}}}}
\newcommand{\NN}[1][]{\ensuremath{\ifstrempty{#1}{\mathbb{N}}{\mathbb{N}^{#1}}}}
\newcommand{\ZZ}[1][]{\ensuremath{\ifstrempty{#1}{\mathbb{Z}}{\mathbb{Z}^{#1}}}}
\newcommand{\QQ}[1][]{\ensuremath{\ifstrempty{#1}{\mathbb{Q}}{\mathbb{Q}^{#1}}}}
\newcommand{\CC}[1][]{\ensuremath{\ifstrempty{#1}{\mathbb{C}}{\mathbb{C}^{#1}}}}
\newcommand{\PP}[1][]{\ensuremath{\ifstrempty{#1}{\mathbb{P}}{\mathbb{P}^{#1}}}}
\newcommand{\HH}[1][]{\ensuremath{\ifstrempty{#1}{\mathbb{H}}{\mathbb{H}^{#1}}}}
\newcommand{\FF}[1][]{\ensuremath{\ifstrempty{#1}{\mathbb{F}}{\mathbb{F}^{#1}}}}
% expected value
\newcommand{\EE}{\ensuremath{\mathbb{E}}}
\newcommand{\charin}{\text{ char }}
\DeclareMathOperator{\sign}{sign}
\DeclareMathOperator{\Aut}{Aut}
\DeclareMathOperator{\Inn}{Inn}
\DeclareMathOperator{\Syl}{Syl}
\DeclareMathOperator{\Gal}{Gal}
\DeclareMathOperator{\GL}{GL} % General linear group
\DeclareMathOperator{\SL}{SL} % Special linear group

%---------------------------------------
% BlackBoard Math Fonts :-
%---------------------------------------

%Captital Letters
\newcommand{\bbA}{\mathbb{A}}	\newcommand{\bbB}{\mathbb{B}}
\newcommand{\bbC}{\mathbb{C}}	\newcommand{\bbD}{\mathbb{D}}
\newcommand{\bbE}{\mathbb{E}}	\newcommand{\bbF}{\mathbb{F}}
\newcommand{\bbG}{\mathbb{G}}	\newcommand{\bbH}{\mathbb{H}}
\newcommand{\bbI}{\mathbb{I}}	\newcommand{\bbJ}{\mathbb{J}}
\newcommand{\bbK}{\mathbb{K}}	\newcommand{\bbL}{\mathbb{L}}
\newcommand{\bbM}{\mathbb{M}}	\newcommand{\bbN}{\mathbb{N}}
\newcommand{\bbO}{\mathbb{O}}	\newcommand{\bbP}{\mathbb{P}}
\newcommand{\bbQ}{\mathbb{Q}}	\newcommand{\bbR}{\mathbb{R}}
\newcommand{\bbS}{\mathbb{S}}	\newcommand{\bbT}{\mathbb{T}}
\newcommand{\bbU}{\mathbb{U}}	\newcommand{\bbV}{\mathbb{V}}
\newcommand{\bbW}{\mathbb{W}}	\newcommand{\bbX}{\mathbb{X}}
\newcommand{\bbY}{\mathbb{Y}}	\newcommand{\bbZ}{\mathbb{Z}}

%---------------------------------------
% MathCal Fonts :-
%---------------------------------------

%Captital Letters
\newcommand{\mcA}{\mathcal{A}}	\newcommand{\mcB}{\mathcal{B}}
\newcommand{\mcC}{\mathcal{C}}	\newcommand{\mcD}{\mathcal{D}}
\newcommand{\mcE}{\mathcal{E}}	\newcommand{\mcF}{\mathcal{F}}
\newcommand{\mcG}{\mathcal{G}}	\newcommand{\mcH}{\mathcal{H}}
\newcommand{\mcI}{\mathcal{I}}	\newcommand{\mcJ}{\mathcal{J}}
\newcommand{\mcK}{\mathcal{K}}	\newcommand{\mcL}{\mathcal{L}}
\newcommand{\mcM}{\mathcal{M}}	\newcommand{\mcN}{\mathcal{N}}
\newcommand{\mcO}{\mathcal{O}}	\newcommand{\mcP}{\mathcal{P}}
\newcommand{\mcQ}{\mathcal{Q}}	\newcommand{\mcR}{\mathcal{R}}
\newcommand{\mcS}{\mathcal{S}}	\newcommand{\mcT}{\mathcal{T}}
\newcommand{\mcU}{\mathcal{U}}	\newcommand{\mcV}{\mathcal{V}}
\newcommand{\mcW}{\mathcal{W}}	\newcommand{\mcX}{\mathcal{X}}
\newcommand{\mcY}{\mathcal{Y}}	\newcommand{\mcZ}{\mathcal{Z}}


%---------------------------------------
% Bold Math Fonts :-
%---------------------------------------

%Captital Letters
\newcommand{\bmA}{\boldsymbol{A}}	\newcommand{\bmB}{\boldsymbol{B}}
\newcommand{\bmC}{\boldsymbol{C}}	\newcommand{\bmD}{\boldsymbol{D}}
\newcommand{\bmE}{\boldsymbol{E}}	\newcommand{\bmF}{\boldsymbol{F}}
\newcommand{\bmG}{\boldsymbol{G}}	\newcommand{\bmH}{\boldsymbol{H}}
\newcommand{\bmI}{\boldsymbol{I}}	\newcommand{\bmJ}{\boldsymbol{J}}
\newcommand{\bmK}{\boldsymbol{K}}	\newcommand{\bmL}{\boldsymbol{L}}
\newcommand{\bmM}{\boldsymbol{M}}	\newcommand{\bmN}{\boldsymbol{N}}
\newcommand{\bmO}{\boldsymbol{O}}	\newcommand{\bmP}{\boldsymbol{P}}
\newcommand{\bmQ}{\boldsymbol{Q}}	\newcommand{\bmR}{\boldsymbol{R}}
\newcommand{\bmS}{\boldsymbol{S}}	\newcommand{\bmT}{\boldsymbol{T}}
\newcommand{\bmU}{\boldsymbol{U}}	\newcommand{\bmV}{\boldsymbol{V}}
\newcommand{\bmW}{\boldsymbol{W}}	\newcommand{\bmX}{\boldsymbol{X}}
\newcommand{\bmY}{\boldsymbol{Y}}	\newcommand{\bmZ}{\boldsymbol{Z}}
%Small Letters
\newcommand{\bma}{\boldsymbol{a}}	\newcommand{\bmb}{\boldsymbol{b}}
\newcommand{\bmc}{\boldsymbol{c}}	\newcommand{\bmd}{\boldsymbol{d}}
\newcommand{\bme}{\boldsymbol{e}}	\newcommand{\bmf}{\boldsymbol{f}}
\newcommand{\bmg}{\boldsymbol{g}}	\newcommand{\bmh}{\boldsymbol{h}}
\newcommand{\bmi}{\boldsymbol{i}}	\newcommand{\bmj}{\boldsymbol{j}}
\newcommand{\bmk}{\boldsymbol{k}}	\newcommand{\bml}{\boldsymbol{l}}
\newcommand{\bmm}{\boldsymbol{m}}	\newcommand{\bmn}{\boldsymbol{n}}
\newcommand{\bmo}{\boldsymbol{o}}	\newcommand{\bmp}{\boldsymbol{p}}
\newcommand{\bmq}{\boldsymbol{q}}	\newcommand{\bmr}{\boldsymbol{r}}
\newcommand{\bms}{\boldsymbol{s}}	\newcommand{\bmt}{\boldsymbol{t}}
\newcommand{\bmu}{\boldsymbol{u}}	\newcommand{\bmv}{\boldsymbol{v}}
\newcommand{\bmw}{\boldsymbol{w}}	\newcommand{\bmx}{\boldsymbol{x}}
\newcommand{\bmy}{\boldsymbol{y}}	\newcommand{\bmz}{\boldsymbol{z}}

%---------------------------------------
% Scr Math Fonts :-
%---------------------------------------

\newcommand{\sA}{{\mathscr{A}}}   \newcommand{\sB}{{\mathscr{B}}}
\newcommand{\sC}{{\mathscr{C}}}   \newcommand{\sD}{{\mathscr{D}}}
\newcommand{\sE}{{\mathscr{E}}}   \newcommand{\sF}{{\mathscr{F}}}
\newcommand{\sG}{{\mathscr{G}}}   \newcommand{\sH}{{\mathscr{H}}}
\newcommand{\sI}{{\mathscr{I}}}   \newcommand{\sJ}{{\mathscr{J}}}
\newcommand{\sK}{{\mathscr{K}}}   \newcommand{\sL}{{\mathscr{L}}}
\newcommand{\sM}{{\mathscr{M}}}   \newcommand{\sN}{{\mathscr{N}}}
\newcommand{\sO}{{\mathscr{O}}}   \newcommand{\sP}{{\mathscr{P}}}
\newcommand{\sQ}{{\mathscr{Q}}}   \newcommand{\sR}{{\mathscr{R}}}
\newcommand{\sS}{{\mathscr{S}}}   \newcommand{\sT}{{\mathscr{T}}}
\newcommand{\sU}{{\mathscr{U}}}   \newcommand{\sV}{{\mathscr{V}}}
\newcommand{\sW}{{\mathscr{W}}}   \newcommand{\sX}{{\mathscr{X}}}
\newcommand{\sY}{{\mathscr{Y}}}   \newcommand{\sZ}{{\mathscr{Z}}}


%---------------------------------------
% Math Fraktur Font
%---------------------------------------

%Captital Letters
\newcommand{\mfA}{\mathfrak{A}}	\newcommand{\mfB}{\mathfrak{B}}
\newcommand{\mfC}{\mathfrak{C}}	\newcommand{\mfD}{\mathfrak{D}}
\newcommand{\mfE}{\mathfrak{E}}	\newcommand{\mfF}{\mathfrak{F}}
\newcommand{\mfG}{\mathfrak{G}}	\newcommand{\mfH}{\mathfrak{H}}
\newcommand{\mfI}{\mathfrak{I}}	\newcommand{\mfJ}{\mathfrak{J}}
\newcommand{\mfK}{\mathfrak{K}}	\newcommand{\mfL}{\mathfrak{L}}
\newcommand{\mfM}{\mathfrak{M}}	\newcommand{\mfN}{\mathfrak{N}}
\newcommand{\mfO}{\mathfrak{O}}	\newcommand{\mfP}{\mathfrak{P}}
\newcommand{\mfQ}{\mathfrak{Q}}	\newcommand{\mfR}{\mathfrak{R}}
\newcommand{\mfS}{\mathfrak{S}}	\newcommand{\mfT}{\mathfrak{T}}
\newcommand{\mfU}{\mathfrak{U}}	\newcommand{\mfV}{\mathfrak{V}}
\newcommand{\mfW}{\mathfrak{W}}	\newcommand{\mfX}{\mathfrak{X}}
\newcommand{\mfY}{\mathfrak{Y}}	\newcommand{\mfZ}{\mathfrak{Z}}
%Small Letters
\newcommand{\mfa}{\mathfrak{a}}	\newcommand{\mfb}{\mathfrak{b}}
\newcommand{\mfc}{\mathfrak{c}}	\newcommand{\mfd}{\mathfrak{d}}
\newcommand{\mfe}{\mathfrak{e}}	\newcommand{\mff}{\mathfrak{f}}
\newcommand{\mfg}{\mathfrak{g}}	\newcommand{\mfh}{\mathfrak{h}}
\newcommand{\mfi}{\mathfrak{i}}	\newcommand{\mfj}{\mathfrak{j}}
\newcommand{\mfk}{\mathfrak{k}}	\newcommand{\mfl}{\mathfrak{l}}
\newcommand{\mfm}{\mathfrak{m}}	\newcommand{\mfn}{\mathfrak{n}}
\newcommand{\mfo}{\mathfrak{o}}	\newcommand{\mfp}{\mathfrak{p}}
\newcommand{\mfq}{\mathfrak{q}}	\newcommand{\mfr}{\mathfrak{r}}
\newcommand{\mfs}{\mathfrak{s}}	\newcommand{\mft}{\mathfrak{t}}
\newcommand{\mfu}{\mathfrak{u}}	\newcommand{\mfv}{\mathfrak{v}}
\newcommand{\mfw}{\mathfrak{w}}	\newcommand{\mfx}{\mathfrak{x}}
\newcommand{\mfy}{\mathfrak{y}}	\newcommand{\mfz}{\mathfrak{z}}


\newcommand{\mytitle}{ECE\_3TC21\_TP Propagation et Antennes}
\newcommand{\myauthor}{github/omgeta}
\newcommand{\mydate}{AU 26-27 S1}

\newcommand{\dd}{\mathrm{d}}
\newcommand{\ve}[1]{\vec{#1}}
\newcommand{\ROS}{\mathrm{ROS}}
\newcommand{\rot}{\operatorname{rot}}
\newcommand{\dv}{\operatorname{div}}

\begin{document}
\raggedright
\raggedcolumns
\footnotesize
\begin{multicols*}{3}
	\setlength{\premulticols}{1pt}
	\setlength{\postmulticols}{1pt}
	\setlength{\multicolsep}{1pt}
	\setlength{\columnsep}{2pt}

	{\normalsize{\textbf{\mytitle}}}\\
	{\footnotesize{\mydate\hspace{2pt}\textemdash\hspace{2pt}\myauthor}}
	%%%%%%%%%%%%%%%%%%%%%%%%%%%%%%%%%%%%%%%%%%%%%%%%%%%%%%
	%                      Begin                         %
	%%%%%%%%%%%%%%%%%%%%%%%%%%%%%%%%%%%%%%%%%%%%%%%%%%%%%%

	\section{Systèmes RF et rappels}
	\textbf{Idée:} un système sans fil émet/reçoit l'information par une onde électromagnétique (OEM) rayonnée/captée par des \textbf{antennes}.
	\begin{enumerate}[\roman*.]
		\item Signaux \textbf{RF}: environ 100\,MHz à 100\,GHz (BF: 1--100\,kHz)
		\item Mobile, Wi-Fi, GPS, Bluetooth: $<6$\,GHz; satellites 1--40\,GHz; militaire 1--100\,GHz
		\item Spectre RF quasi entièrement attribué: \textbf{saturé}
	\end{enumerate}
	\begin{gather*}
		\lambda=\frac{c}{f}=\frac{\lambda_0}{\sqrt{\varepsilon_r}},\quad c=3\cdot10^8\text{\,m/s}
	\end{gather*}
	\begin{tabular}{|l|l||l|l|}
		\hline
		$f$     & $\lambda_0$ & $f$     & $\lambda_0$ \\
		\hline
		1\,kHz  & 300\,km     & 2,5\,GHz & 12\,cm      \\
		100\,MHz & 3\,m       & 10\,GHz & 3\,cm       \\
		1\,GHz  & 30\,cm      & 100\,GHz & 3\,mm      \\
		\hline
	\end{tabular}

	\textbf{Pourquoi la RF?}
	\begin{enumerate}[\roman*.]
		\item Taille d'antenne $\propto\lambda\propto1/f$ ($\sim\lambda/2$): petite, intégrable
		\item Propagation correcte (murs, vitres, végétation), mais affaiblissement $\nearrow$ avec $f$
		\item Grandes bandes disponibles loin de 0\,Hz: \textbf{haut débit}
	\end{enumerate}

	\textbf{Spécificité RF}: dimensions des pistes, câbles, composants $\sim\lambda$ ou plus $\Rightarrow$ théorie des \textbf{lignes} (§2), facteur de réflexion, adaptation (§3), paramètres S (§4).
	\begin{enumerate}[\roman*.]
		\item On mesure des \textbf{puissances} (analyseur de spectre, wattmètre), pas $V$/$I$
		\item \textbf{VNA} (analyseur de réseaux vectoriel): mesure les paramètres S; étalonnage obligatoire avant chaque usage
	\end{enumerate}

	\colbreak
	\textbf{Affaiblissement}: le problème fondamental; remèdes: \textbf{amplification} et \textbf{adaptation}.
	\begin{enumerate}[\roman*.]
		\item Câble: pertes (Joule, effet de peau, $\varepsilon$) $\nearrow$ avec $f$: 0,05\,dB/m à 10\,MHz, 0,64 à 1\,GHz, 1,6 à 6\,GHz, 6,4 à 67\,GHz
		\item Espace libre (FSPL) bien pire, en $f^2d^2$ (§6):
	\end{enumerate}
	\begin{gather*}
		\mathrm{FSPL}_{dB}=20\log\frac{4\pi df}{c}=20\log d+20\log f-147{,}5
	\end{gather*}
	\textbf{Exemple}: $d=100$\,m: 52,5\,dB à 100\,MHz, 72,5\,dB à 1\,GHz (100$\times$ plus atténué)

	\subsection{Chaîne d'émission/réception}
	\begin{enumerate}[\roman*.]
		\item \textbf{TX}: bits $\to$ MOD $\to$ DAC (signal BF) $\to$ mélangeur (+ oscillateur local LO, transpose vers RF) $\to$ \textbf{PA} (amplificateur de puissance) $\to$ antenne
		\item \textbf{RX}: antenne $\to$ \textbf{LNA} (faible bruit) $\to$ mélangeur (RF$\to$BF) $\to$ ADC $\to$ DEMOD $\to$ bits
	\end{enumerate}

	\subsection{Décibels}
	Rapport de puissances: $G_{dB}=10\log\frac{P_2}{P_1}$; puissance absolue:
	\begin{gather*}
		P_{dBm}=10\log\frac{P}{1\,\text{mW}},\quad P_{dBW}=10\log\frac{P}{1\,\text{W}}\\
		P_{dBm}=P_{dBW}+30
	\end{gather*}
	\begin{tabular}{|l|l||l|l|}
		\hline
		$\times2$  & $+3$\,dB  & 1\,mW   & 0\,dBm, $-30$\,dBW \\
		$\times4$  & $+6$\,dB  & 100\,mW & 20\,dBm            \\
		$\times10$ & $+10$\,dB & 1\,W    & 30\,dBm, 0\,dBW    \\
		$\times100$ & $+20$\,dB & 0,5\,W & 27\,dBm           \\
		\hline
	\end{tabular}
	\begin{enumerate}[\roman*.]
		\item dB$\pm$dB = dB; dBm$\pm$dB = dBm; dBm$-$dBm = dB; jamais dBm$+$dBm
		\item Tension/champ: $20\log$ (ex. $|\Gamma|_{dB}=20\log|\Gamma|$)
		\item Wi-Fi: TX max 20\,dBm; RX $-30$ (excellent) à $-90$\,dBm
	\end{enumerate}

	\subsection{Méthode}
	\begin{enumerate}[\roman*.]
		\item $\lambda$ dans le bon milieu: $\lambda_0/\sqrt{\varepsilon_r}$
		\item Produits de gains/pertes: passer en dB/dBm
	\end{enumerate}

	\colbreak
	\subsection{Rappels OEM (hors contrôle)}
	Milieu \textbf{MLHI}: linéaire ($\varepsilon,\mu$ indépendants de l'amplitude), homogène, isotrope.
	\begin{gather*}
		\ve D=\varepsilon\ve E,\quad\ve B=\mu\ve H,\quad\ve j=\sigma\ve E\\
		\varepsilon=\varepsilon_0\varepsilon_r,\ \varepsilon_0=8{,}854\cdot10^{-12}\text{\,F/m};\ \mu=\mu_0=4\pi\cdot10^{-7}\text{\,H/m}
	\end{gather*}
	\begin{enumerate}[\roman*.]
		\item Pertes diélectriques: $\varepsilon=\varepsilon'-j\varepsilon''$, $\tan\delta=\varepsilon''/\varepsilon'$ (Téflon $5\cdot10^{-4}$, air 0)
		\item Optique: $n=\sqrt{\varepsilon_r}$; diélectrique: $\mu_r=1$
	\end{enumerate}

	\textbf{Maxwell} (diélectrique parfait: $\rho=0$, $\ve j=0$):
	\begin{gather*}
		\rot\ve E=-\mu_0\frac{\partial\ve H}{\partial t},\quad\rot\ve H=\varepsilon\frac{\partial\ve E}{\partial t}\\
		\dv\ve E=0,\quad\dv\ve H=0
	\end{gather*}
	\textbf{OPPM} (onde plane progressive monochromatique) selon $Oz$:
	\begin{gather*}
		\ve E(z,t)=\ve E_+e^{j(\omega t-kz)},\quad\ve H=\frac{1}{\eta}\,\ve u_z\wedge\ve E\\
		\Delta\ve E+\omega^2\mu_0\varepsilon\ve E=0\ \text{(Helmholtz)}\\
		k=\omega\sqrt{\mu_0\varepsilon_0\varepsilon_r}=\frac{2\pi}{\lambda},\quad v_\varphi=\frac{\omega}{k}=\frac{c}{\sqrt{\varepsilon_r}}
	\end{gather*}
	\begin{enumerate}[\roman*.]
		\item \textbf{TEM}: $\ve E\perp\ve H\perp\ve k$, trièdre direct, pas de composante longitudinale
		\item Vitesse de groupe $v_g=\dd\omega/\dd k$ (énergie, information); non dispersif: $v_g=v_\varphi$; dispersif: $v_g<v_\varphi$, $v_gv_\varphi=c^2/\varepsilon_r$
		\item Impédance d'onde: $\eta=\frac{|E|}{|H|}=\sqrt{\frac{\mu}{\varepsilon}}=\frac{\eta_0}{\sqrt{\varepsilon_r}}$, $\eta_0=377\,\Omega$
	\end{enumerate}
	\textbf{Poynting} (densité surfacique de puissance, W/m²):
	\begin{gather*}
		\langle\ve\Pi\rangle=\tfrac12\Re(\ve E\wedge\ve H^*),\quad P=\frac{E_+^2}{2\eta}S
	\end{gather*}
	\textbf{Zones}: champ proche (Rayleigh, Fresnel) puis \textbf{champ lointain} (Fraunhofer) au-delà de $2D^2/\lambda$ ($D$: plus grande dimension de l'antenne): onde sphérique en $1/r$, localement plane.


	\colbreak
	\section{Lignes de transmission}
	\textbf{Idée:} dès qu'une connexion mesure $\sim\lambda$, tension et courant se \textbf{propagent} (ondes) le long de la ligne.
	\begin{enumerate}[\roman*.]
		\item 1 conducteur: guide d'onde ($>30$\,GHz, radar, satellite; pertes min., forte puissance)
		\item 2 conducteurs + diélectrique $\varepsilon_r$ ($\mu_r=1$): \textbf{coaxial} (TEM, blindé) et \textbf{microstrip} (quasi-TEM, champ en partie dans l'air: rayonne)
	\end{enumerate}

	\subsection{Modèle RLGC}
	Tronçon $\dd z\to0$, éléments \textbf{distribués} (par mètre): $R$ ($\Omega$/m) pertes métalliques; $L$ (H/m) champ magnétique du courant; $C$ (F/m) champ électrique entre conducteurs; $G$ (S/m) pertes diélectriques. \textbf{Sans pertes}: $R=G=0$.

	Coaxial (rayons $a$ intérieur, $b$ extérieur):
	\begin{gather*}
		L=\frac{\mu}{2\pi}\ln\frac ba,\quad C=\frac{2\pi\varepsilon}{\ln(b/a)}
	\end{gather*}

	\subsection{Équations des télégraphistes}
	Kirchhoff sur un tronçon, régime harmonique:
	\begin{gather*}
		-\frac{\dd V}{\dd z}=(R+j\omega L)I,\quad-\frac{\dd I}{\dd z}=(G+j\omega C)V\\
		\frac{\dd^2V}{\dd z^2}-\gamma^2V=0,\quad\frac{\dd^2I}{\dd z^2}-\gamma^2I=0\\
		\gamma=\sqrt{(R+j\omega L)(G+j\omega C)}=\alpha+j\beta
	\end{gather*}
	$\alpha$ (Np/m): \textbf{atténuation}; $\beta$ (rad/m): \textbf{déphasage}. Sans pertes: $\gamma=j\beta$,
	\begin{gather*}
		\beta=\omega\sqrt{LC}=\frac{2\pi}{\lambda},\quad v_\varphi=\frac1{\sqrt{LC}}=\frac{c}{\sqrt{\varepsilon_r}}
	\end{gather*}
	\textbf{Solutions}: onde \textbf{incidente} (vers la charge, $z>0$) + onde \textbf{réfléchie}:
	\begin{gather*}
		V(z)=V_0^+e^{-\gamma z}+V_0^-e^{+\gamma z}\\
		I(z)=\frac{V_0^+}{Z_C}e^{-\gamma z}-\frac{V_0^-}{Z_C}e^{+\gamma z}
	\end{gather*}

	\colbreak
	\subsection{Impédance caractéristique}
	\begin{gather*}
		Z_C=\frac{V^+(z)}{I^+(z)}=-\frac{V^-(z)}{I^-(z)}=\sqrt{\frac{R+j\omega L}{G+j\omega C}}\\
		\text{sans pertes: }Z_C=\sqrt{\frac LC}\ \text{(réelle)},\quad\text{coaxial: }Z_C=\frac{60}{\sqrt{\varepsilon_r}}\ln\frac ba
	\end{gather*}
	\begin{enumerate}[\roman*.]
		\item Constante le long de la ligne; pas «matérielle»: répartition de la puissance de l'onde progressive entre $V^+$ et $I^+$; impédance vue d'une ligne \textbf{infinie}
		\item \textbf{50\,$\Omega$} standard: compromis pertes min. (77\,$\Omega$) / puissance max. (30\,$\Omega$). Téflon, $a=1{,}4$, $b=4{,}9$\,mm: $\approx52\,\Omega$
	\end{enumerate}

	\subsection{Facteur de réflexion}
	Convention: charge $Z_L$ en $z=0$, générateur côté $z<0$; on se place à $z=-l$.
	\begin{gather*}
		\Gamma_L=\Gamma(0)=\frac{V_0^-}{V_0^+}=\frac{Z_L-Z_C}{Z_L+Z_C},\quad Z_L=Z_C\frac{1+\Gamma_L}{1-\Gamma_L}\\
		\Gamma(z)=\Gamma_Le^{+2j\beta z}\Rightarrow\Gamma(-l)=\Gamma_Le^{-2j\beta l}
	\end{gather*}
	\begin{enumerate}[\roman*.]
		\item CC ($Z_L=0$): $\Gamma=-1$; CO ($Z_L=\infty$): $\Gamma=+1$; adapté ($Z_L=Z_C$): $\Gamma=0$
		\item $|\Gamma|\leq1$ (charge passive); sans pertes, $|\Gamma|$ constant le long de la ligne, seule la phase tourne
	\end{enumerate}

	\subsection{Puissance (ligne sans pertes)}
	\begin{gather*}
		P(z)=\tfrac12\Re[V(z)I^*(z)]=\frac{|V_0^+|^2}{2Z_C}\big(1-|\Gamma|^2\big)\\
		P_L=P_{inc}\big(1-|\Gamma_L|^2\big),\quad P_{ref}=|\Gamma_L|^2P_{inc}
	\end{gather*}
	$P_{inc}=|V_0^+|^2/2Z_C$: puissance incidente (celle du générateur s'il est adapté).

	\begin{tabular}{|l|l|l|l|}
		\hline
		$|\Gamma|_{dB}$ & $|\Gamma|$ & ROS  & $P_{ref}$ \\
		\hline
		$-6$          & 0,50       & 3    & 25\,\%     \\
		$-9{,}5$      & 0,33       & 2    & 11\,\%     \\
		$-10$         & 0,32       & 1,92 & 10\,\%     \\
		$-15{,}6$     & 0,17       & 1,4  & 2,8\,\%    \\
		$-20$         & 0,10       & 1,22 & 1\,\%      \\
		\hline
	\end{tabular}

	\colbreak
	\subsection{Impédance ramenée}
	Impédance vue à la distance $l$ de la charge (sans pertes):
	\begin{gather*}
		Z(-l)=\frac{V(-l)}{I(-l)}=Z_C\frac{Z_L+jZ_C\tan\beta l}{Z_C+jZ_L\tan\beta l}
	\end{gather*}
	\begin{tabular}{|l|l|}
		\hline
		\textbf{Cas}             & $Z(-l)$                    \\
		\hline
		$Z_L=Z_C$                & $Z_C$ (tout $l$)           \\
		$Z_L=0$ (CC)             & $jZ_C\tan\beta l$          \\
		$Z_L=\infty$ (CO)        & $-jZ_C\cot\beta l$         \\
		$l=\lambda/2$            & $Z_L$ (périodicité)        \\
		$l=\lambda/4$            & $Z_C^2/Z_L$ (transformateur) \\
		\hline
	\end{tabular}
	\begin{enumerate}[\roman*.]
		\item $\lambda/4$: CC$\leftrightarrow$CO; capacité $C\to$ inductance $L_{eq}=Z_C^2C$; inductance $L\to$ capacité $C_{eq}=L/Z_C^2$
		\item Tronçon CC/CO: réactance pure réglable par $l$ (stub, §3)
		\item Démonstration: $V$, $I$ en $-l$, factoriser $V_0^+$, remplacer $\Gamma_L$, puis $e^{\pm j\beta l}\to\cos,\sin$
	\end{enumerate}

	\subsection{Ondes stationnaires}
	Si $Z_L\neq Z_C$: incidente + réfléchie $\Rightarrow$ \textbf{onde stationnaire}: $|V(z)|$ présente des ventres ($V^+$, $V^-$ en phase: $V_{max}=|V^+|+|V^-|$) et des nœuds (opposition: $V_{min}=|V^+|-|V^-|$) fixes; ventre$\to$ventre: $\lambda/2$, ventre$\to$nœud: $\lambda/4$.
	\begin{gather*}
		\ROS=\frac{V_{max}}{V_{min}}=\frac{1+|\Gamma|}{1-|\Gamma|},\quad|\Gamma|=\frac{\ROS-1}{\ROS+1}
	\end{gather*}
	$1\leq\ROS<\infty$; adapté: 1; CC/CO: $\infty$. Anglais: VSWR.

	\subsection{Méthode}
	\begin{enumerate}[\roman*.]
		\item $\lambda=\lambda_0/\sqrt{\varepsilon_r}$, $\beta=2\pi/\lambda$; $l$ en fraction de $\lambda$
		\item $\Gamma_L=\frac{Z_L-Z_C}{Z_L+Z_C}$, puis ROS, $|\Gamma|_{dB}$, $P_{ref}=|\Gamma|^2P_{inc}$
		\item Impédance ailleurs: $Z(-l)$ (cas $\lambda/4$, $\lambda/2$, CC, CO d'abord)
		\item Onde progressive seule: $Z(z)=Z_C$ partout, toute la puissance va dans la charge
	\end{enumerate}

	\colbreak
	\section{Adaptation d'impédance}
	\textbf{But:} transmettre le maximum de la puissance du générateur ($Z_G$) à la charge $Z_L$ (ampli, filtre, antenne\dots): supprimer l'onde réfléchie (puissance perdue, risque de destruction de composants, batterie).
	\begin{enumerate}[\roman*.]
		\item Hypothèse du cours: générateur adapté $Z_G=Z_C$ (50\,$\Omega$)
		\item On insère un quadripôle $Q_{match}$ \textbf{sans pertes} (éléments réactifs: $L$, $C$ ou tronçons de ligne) tel que l'impédance vue $Z_E=Z_C$, soit $\Gamma_E=0$
		\item $Z_L$ doit avoir une \textbf{partie réelle}: une charge purement réactive ne peut absorber de puissance, donc pas d'adaptation possible
		\item Solutions à \textbf{bande étroite} (valables autour d'une $f$)
		\item BF: «adaptation» = $Z_{in}$ grande pour ne pas écraser une tension; RF: transfert de puissance
	\end{enumerate}

	\subsection{Ligne quart d'onde}
	Ligne $\lambda/4$ d'impédance $Z'_C$ entre la ligne et $Z_L$:
	\begin{gather*}
		Z_E=\frac{Z_C'^2}{Z_L}=Z_C\ \Rightarrow\ Z'_C=\sqrt{Z_LZ_C}
	\end{gather*}
	\begin{enumerate}[\roman*.]
		\item $Z_L$ \textbf{réelle} uniquement ($Z'_C$ réelle)
		\item Longueur: $l=\lambda/4$ avec $\lambda$ \textbf{dans le matériau} ($\varepsilon_r$)
		\item $Z_L$ complexe: d'abord annuler sa partie imaginaire (réactance $-jX$ en série, ou tronçon de ligne $Z_C$ de longueur $l'$ qui la rend réelle), puis $\lambda/4$
	\end{enumerate}
	\textbf{Exemple} (TD1): $Z_L=100\,\Omega$, 2\,GHz, $\varepsilon_r=2{,}1$:
	\begin{enumerate}[\roman*.]
		\item Sans adaptation: $\Gamma_L=\frac13$, ROS$=2$, 11\,\% réfléchi (88,9\,mW reçus sur 20\,dBm)
		\item $Z'_C=\sqrt{100\cdot50}=70{,}7\,\Omega$
		\item $\lambda=\frac{0{,}15}{\sqrt{2{,}1}}=10{,}35$\,cm, $l=25{,}9$\,mm
	\end{enumerate}

	\colbreak
	\textbf{Exemple} (TD1): antenne $R=50\,\Omega\parallel C=10$\,pF à 2\,GHz.
	\begin{enumerate}[\roman*.]
		\item $Y=\frac1R+jC\omega$: ajouter $L\parallel$ avec $L\omega=\frac1{C\omega}$ annule la susceptance $\Rightarrow Y=1/50$
		\item $L=\frac1{\omega^2C}=0{,}63$\,nH
		\item Avec un condensateur: $C'$ au bout d'une ligne $\lambda/4$ (50\,$\Omega$) $\equiv$ inductance $Z_C^2C'$: $C'=L/Z_C^2=0{,}25$\,pF
	\end{enumerate}

	\subsection{Stub}
	\textbf{Stub}: tronçon de ligne sans pertes (même $Z_C$, $\beta$) en \textbf{parallèle} sur la ligne, terminé par CC ou CO, à la distance $d$ de la charge, de longueur $d'$.
	\begin{enumerate}[\roman*.]
		\item Éléments en parallèle: raisonner en \textbf{admittance}
		\item Condition: $Y_E=Y_{Z_L}(-d)+Y_{stub}(d')=Y_C=\frac1{Z_C}$
		\item $d$ annule $\Re$: $\Re\,Y_{Z_L}(-d)=Y_C$; $d'$ annule $\Im$
	\end{enumerate}
	\begin{gather*}
		Y_{Z_L}(-d)=\frac1{Z_C}\,\frac{Z_C+jZ_L\tan\beta d}{Z_L+jZ_C\tan\beta d}\\
		\text{CC: }Y_{stub}=\frac{1}{jZ_C\tan\beta d'},\quad\text{CO: }Y_{stub}=\frac{j\tan\beta d'}{Z_C}
	\end{gather*}
	\textbf{$Z_L$ réelle}, $r=Z_L/Z_C$, $t=\tan\beta d$:
	\begin{gather*}
		y=Z_CY_{Z_L}=\frac{r(1+t^2)+jt(r^2-1)}{r^2+t^2}\\
		\Re y=1\Rightarrow\tan\beta d=\pm\sqrt r\\
		\Im y=b=\pm\frac{r-1}{\sqrt r}
	\end{gather*}
	Annuler $b$ avec le stub:
	\begin{gather*}
		\text{CC: }\tan\beta d'=\frac1b=\pm\frac{\sqrt r}{r-1}\\
		\text{CO: }\tan\beta d'=-b=\mp\frac{r-1}{\sqrt r}
	\end{gather*}
	\begin{enumerate}[\roman*.]
		\item Même signe $\pm$ que pour $d$
		\item $d$, $d'$ définis modulo $\lambda/2$: prendre l'angle dans $[0,\pi[$, puis $l=\frac{\text{angle}}{2\pi}\lambda$
		\item On garde $d$, $d'$ courts (lignes imparfaites); choix CC/CO selon les longueurs obtenues (donné en contrôle)
	\end{enumerate}

	\colbreak
	\subsection{Éléments localisés (réseau en L)}
	Deux réactances pures $jX_S$ (série) et $jX_P$ (parallèle), pas de résistance (pas de dissipation):
	\begin{enumerate}[\roman*.]
		\item $Z_L>Z_C$: $jX_P$ \textbf{parallèle} sur $Z_L$, puis $jX_S$ en série
		\item $Z_L<Z_C$: $jX_S$ \textbf{série} avec $Z_L$, puis $jX_P$ en parallèle
		\item $X_S$ et $X_P$ de signes opposés: 2 solutions ($L$ série/$C$ parallèle ou l'inverse)
		\item Choix des composants selon les valeurs réalisables
	\end{enumerate}
	\textbf{$Z_L=R_L>Z_E$} (on veut $Z_E$, ici 50\,$\Omega$):
	\begin{gather*}
		X_P=\pm R_L\sqrt{\frac{Z_E}{R_L-Z_E}},\quad X_S=\mp\sqrt{Z_E(R_L-Z_E)}
	\end{gather*}
	\textbf{$Z_L=R_L<Z_E$}:
	\begin{gather*}
		X_S=\pm\sqrt{R_L(Z_E-R_L)},\quad X_P=\mp Z_E\sqrt{\frac{R_L}{Z_E-R_L}}
	\end{gather*}
	Méthode: écrire $Z_E=jX_S+\frac{R_L\,jX_P}{R_L+jX_P}$ (ou $\frac1{Z_E}=\frac1{jX_P}+\frac1{R_L+jX_S}$), séparer réel et imaginaire.
	\begin{gather*}
		X>0:\ L=\frac{X}{\omega};\qquad X<0:\ C=\frac{1}{\omega|X|}
	\end{gather*}
	\textbf{Exemple} (1\,GHz, $\omega=6{,}28\cdot10^9$):

	\begin{tabular}{|l|l|l|}
		\hline
		                & $Z_L=200\,\Omega$     & $Z_L=35\,\Omega$     \\
		\hline
		$|X_P|$         & 115,5\,$\Omega$       & 76,4\,$\Omega$       \\
		$|X_S|$         & 86,6\,$\Omega$        & 22,9\,$\Omega$       \\
		\hline
		Sol. 1          & $L_P=18{,}4$\,nH      & $C_P=2{,}08$\,pF     \\
		                & $C_S=1{,}84$\,pF      & $L_S=3{,}65$\,nH     \\
		\hline
		Sol. 2          & $C_P=1{,}38$\,pF      & $L_P=12{,}2$\,nH     \\
		                & $L_S=13{,}8$\,nH      & $C_S=6{,}95$\,pF     \\
		\hline
	\end{tabular}

	\subsection{Méthode}
	\begin{enumerate}[\roman*.]
		\item $Z_L$ réelle? $\lambda/4$ direct; sinon annuler $\Im$ d'abord
		\item Stub: admittances, $\Re$ par $d$, $\Im$ par $d'$
		\item Réseau en L: comparer $R_L$ à 50\,$\Omega$ pour placer $X_P$ côté grande résistance
		\item Vérifier: recalculer $Z_E$, doit valoir 50\,$\Omega$
	\end{enumerate}

	\colbreak
	\section{Paramètres S}
	\textbf{Idée:} en hyperfréquence, CC/CO larges bandes, sources de courant et mesures de tension sont impossibles: on caractérise un multipôle par des \textbf{ondes de puissance} entrantes/sortantes (mesurables).
	\begin{enumerate}[\roman*.]
		\item Matrice $Z$ (BF): $Z_{11}=V_1/i_1$ avec accès 2 en CO, donc inutilisable en RF
		\item Paramètres S: accès fermés sur l'\textbf{impédance de référence} $Z_0$ (réelle, en général $Z_0=Z_C=50\,\Omega$)
	\end{enumerate}

	\subsection{Ondes de puissance}
	À l'accès $k$: $V_k=V_k^++V_k^-$, $I_k=\frac{1}{Z_0}(V_k^+-V_k^-)$.
	\begin{gather*}
		a_k=\frac{V_k+Z_0I_k}{2\sqrt{Z_0}}=\frac{V_k^+}{\sqrt{Z_0}},\quad b_k=\frac{V_k-Z_0I_k}{2\sqrt{Z_0}}=\frac{V_k^-}{\sqrt{Z_0}}\\
		P_k^+=\frac{|V_k^+|^2}{2Z_0}=\tfrac12|a_k|^2,\quad P_k^-=\tfrac12|b_k|^2
	\end{gather*}
	$a$, $b$ complexes: gardent la phase, donc dépendent du \textbf{plan de référence}.

	\subsection{Matrice S}
	\begin{gather*}
		(b)=S\,(a):\quad b_i=\sum_jS_{ij}a_j\\
		\begin{cases}
			b_1=S_{11}a_1+S_{12}a_2 \\
			b_2=S_{21}a_1+S_{22}a_2
		\end{cases}
	\end{gather*}
	\begin{enumerate}[\roman*.]
		\item Complexes, dépendent de $f$, du plan de référence et de $Z_0$
		\item Multipôle linéaire, régime harmonique; $N$ accès: matrice $N\times N$
	\end{enumerate}
	\textbf{Dipôle} (1 accès, ex. antenne):
	\begin{gather*}
		S_{11}=\frac{b_1}{a_1}=\frac{V_1^-}{V_1^+}=\Gamma_1
	\end{gather*}

	\colbreak
	\textbf{Quadripôle}: alimenter un accès, fermer l'autre sur $Z_0$ (pas de réflexion: onde entrante nulle):
	\begin{gather*}
		S_{11}=\frac{b_1}{a_1}\Big|_{a_2=0}=\Gamma_1,\quad S_{21}=\frac{b_2}{a_1}\Big|_{a_2=0}=T_{21}\\
		S_{22}=\frac{b_2}{a_2}\Big|_{a_1=0}=\Gamma_2,\quad S_{12}=\frac{b_1}{a_2}\Big|_{a_1=0}=T_{12}
	\end{gather*}
	\begin{enumerate}[\roman*.]
		\item $S_{ii}$: \textbf{réflexion} à l'accès $i$
		\item $S_{21}$: \textbf{transmission directe} (1$\to$2); $S_{12}$: inverse
		\item Mesure au \textbf{VNA}: générateur adapté sur un accès, autres accès sur charges adaptées
	\end{enumerate}

	\subsection{Propriétés}
	\begin{tabular}{|l|l|}
		\hline
		\textbf{Adapté} (accès $i$) & $S_{ii}=0$ (pratique: $<-10$\,dB) \\
		\textbf{Réciproque}       & $S_{ij}=S_{ji}$ ($S=S^T$)                   \\
		\textbf{Symétrique}       & réciproque et $S_{ii}=S_{jj}$               \\
		\textbf{Passif}           & tous $|S_{ij}|\leq1$                        \\
		\textbf{Actif}            & au moins un $|S_{ij}|>1$                    \\
		\textbf{Sans pertes}      & $\sum_i|S_{ij}|^2=1$ (colonnes)             \\
		\hline
	\end{tabular}
	\begin{enumerate}[\roman*.]
		\item Non réciproque: ampli, isolateur, circulateur; unilatéral: $S_{12}=0$
	\end{enumerate}

	\subsection{Puissances}
	Accès 2 fermé sur $Z_0$, $P_{inc1}=\frac12|a_1|^2$:
	\begin{gather*}
		P_{entr\acute ee}=P_{inc1}(1-|S_{11}|^2)\\
		P_{charge}=|S_{21}|^2P_{inc1},\quad G=|S_{21}|^2\\
		G_{dB}=20\log|S_{21}|,\quad|S_{11}|_{dB}=20\log|S_{11}|
	\end{gather*}
	Pertes d'insertion $=-20\log|S_{21}|$; pertes par retour (return loss) $=-20\log|S_{11}|$.

	\colbreak
	\subsection{Exemples (TD3)}
	\textbf{Ligne} sans pertes, longueur $l$, $Z_C=Z_0=50\,\Omega$:
	\begin{enumerate}[\roman*.]
		\item Fermée sur $Z_0$: $Z_E=Z_C$, $\Gamma_E=0\Rightarrow S_{11}=0$
		\item $S_{21}=\frac{V_2}{V_1}$ (car $a_2=0$, $b_1=0$): $S_{21}=e^{-j\beta l}$, simple retard de phase
		\item Symétrique: $S_{22}=0$, $S_{12}=S_{21}$
	\end{enumerate}
	\begin{gather*}
		S_{ligne}=\begin{pmatrix}0&e^{-j\beta l}\\e^{-j\beta l}&0\end{pmatrix}
	\end{gather*}
	\begin{enumerate}[\roman*.]
		\item Si $Z_0\neq Z_C$: $S_{11},S_{22}\neq0$, $|S_{21}|<1$
	\end{enumerate}

	\textbf{Dispositif T}: $S=\begin{pmatrix}0{,}6&0\\14e^{j135^\circ}&0\end{pmatrix}$
	\begin{enumerate}[\roman*.]
		\item 2 accès; $S_{12}\neq S_{21}$: non réciproque; $|S_{21}|=14>1$: actif
		\item $|S_{11}|_{dB}=20\log0{,}6=-4{,}4$\,dB $>-10$\,dB: pas adapté en entrée
		\item $S_{12}=0$, gain: \textbf{amplificateur}; $G=14^2=196=22{,}9$\,dB
		\item $P_{entr\acute ee}=(1-0{,}36)P_{inc1}=0{,}64\,P_{inc1}$
	\end{enumerate}

	\textbf{Circulateur} $S=\begin{pmatrix}0&0&1\\1&0&0\\0&1&0\end{pmatrix}$:
	\begin{enumerate}[\roman*.]
		\item 3 accès; $S_{ii}=0$: adapté; $S_{21}\neq S_{12}$: non réciproque; $|S_{ij}|\leq1$: passif (et sans pertes)
		\item $b_2=a_1$, $b_3=a_2$, $b_1=a_3$: $1\to2\to3\to1$, l'onde ne tourne que dans un sens
		\item \textbf{Radar monostatique} (une antenne TX/RX): émetteur sur 1, antenne sur 2, récepteur sur 3; l'écho entre en 2, sort en 3, isolé de l'émetteur
	\end{enumerate}

	\subsection{Méthode}
	\begin{enumerate}[\roman*.]
		\item Lire $|S|$: linéaire ou dB ($20\log$)?
		\item Diagonale $\to$ adaptation; $S$ vs $S^T$ $\to$ réciprocité; $|S|>1$ $\to$ actif
		\item Fonction: ampli ($|S_{21}|>1$), atténuateur ($|S_{21}|<1$), isolateur, circulateur
		\item Puissance: $\times(1-|S_{11}|^2)$ entrée, $\times|S_{21}|^2$ charge
	\end{enumerate}

	\colbreak
	\section{Antennes}
	\textbf{Antenne}: convertit une onde \textbf{guidée} (coaxial, microstrip, guide d'onde) en onde \textbf{rayonnée}, et inversement.
	\begin{enumerate}[\roman*.]
		\item Rayonnement: charges \textbf{accélérées}, seules les sources variables dans le temps rayonnent
		\item Contraintes: bande(s) de fréquence, taille ($\sim\lambda/2$: 43\,cm à 700\,MHz), environnement (température, puissance)
	\end{enumerate}
	Trois caractéristiques à l'émission:
	\begin{enumerate}[\roman*.]
		\item \textbf{Impédance d'entrée}: adaptation au générateur
		\item \textbf{Diagramme de rayonnement}: répartition spatiale de la puissance
		\item \textbf{Polarisation}: orientation du champ $\ve E$
	\end{enumerate}

	\subsection{Impédance d'entrée, bande passante}
	\begin{gather*}
		Z_{ant}=R_{ant}+jX_{ant},\quad\Gamma_{ant}=\frac{Z_{ant}-Z_C}{Z_{ant}+Z_C}=S_{11}\\
		P_{ant}=(1-|\Gamma_{ant}|^2)P_{gen},\quad P_{ray}=e_{ray}P_{ant}
	\end{gather*}
	\begin{enumerate}[\roman*.]
		\item $e_{ray}\in\,]0,1[$: efficacité de rayonnement (pertes métalliques, diélectriques)
		\item $Z_{ant}$ dépend de $f$ et de l'environnement: $\Gamma=0$ impossible sur toute la bande
		\item \textbf{Bande passante} (en adaptation): $|\Gamma_{ant}|_{dB}<-10$\,dB (mobile: $-6$\,dB), ou $\ROS<2$ ($\Leftrightarrow-9{,}5$\,dB)
	\end{enumerate}
	\textbf{Exemple} (autonomie TD4): cahier des charges 3,1--5\,GHz:
	\begin{enumerate}[\roman*.]
		\item ROS$<2$: lire la plage où $|\Gamma|<-9{,}5$\,dB, doit contenir 3,1--5\,GHz
		\item ROS$<1{,}4$: seuil $|\Gamma|=\frac{0{,}4}{2{,}4}=0{,}17=-15{,}6$\,dB, bande plus étroite
	\end{enumerate}

	\colbreak
	\subsection{Diagramme de rayonnement}
	Tracé en \textbf{champ lointain} ($r>2D^2/\lambda$), coordonnées sphériques $(r,\theta,\varphi)$.
	\begin{gather*}
		p(r,\theta,\varphi)=\|\langle\ve\Pi\rangle\|=\tfrac12\|\Re(\ve E\wedge\ve H^*)\|\\
		P_{ray}=\int_0^{2\pi}\!\!\int_0^\pi p\,r^2\sin\theta\,\dd\theta\,\dd\varphi,\quad p_{iso}=\frac{P_{ray}}{4\pi r^2}
	\end{gather*}
	\textbf{Directivité}, \textbf{gain} (référence: antenne \textbf{isotrope}, uniforme dans tout l'espace):
	\begin{gather*}
		D(\theta,\varphi)=\frac{p}{P_{ray}/4\pi r^2},\quad G(\theta,\varphi)=\frac{p}{P_{ant}/4\pi r^2}=e_{ray}D\\
		D_{dBi}=10\log D
	\end{gather*}
	\begin{enumerate}[\roman*.]
		\item $D$ ne peut pas être $>1$ partout: concentrer dans une direction se paie ailleurs
		\item Représentations: 3D (qualitatif), coupes polaires ou cartésiennes; normalisé: max à 0\,dB
		\item \textbf{Lobe principal} entre les deux minima autour du max; \textbf{lobes secondaires} de part et d'autre (niveau absolu, ou relatif au max)
		\item \textbf{Ouverture à $-3$\,dB} $\theta_{3dB}$: plage où $G\geq G_{max}-3$\,dB (mi-puissance)
		\item \textbf{Rapport avant/arrière}: $G_{avant}-G_{arri\grave ere}$ (dB)
		\item Omnidirectionnel (max identique dans un plan) $\neq$ isotrope
	\end{enumerate}

	\colbreak
	\subsection{Polarisation}
	Direction de $\ve E$; propagation selon $Oz$:
	\begin{gather*}
		\ve E=E_{0x}e^{j(\omega t-kz)}\ve u_x+E_{0y}e^{j(\omega t-kz+\varphi)}\ve u_y
	\end{gather*}
	\begin{enumerate}[\roman*.]
		\item \textbf{Rectiligne}: $\varphi=0$ ou $\pi$ (ou une composante nulle): l'extrémité de $\ve E$ décrit une droite
		\item \textbf{Circulaire}: $\varphi=\pm\pi/2$ et $E_{0x}=E_{0y}$: un cercle
		      \begin{itemize}
			      \item $\varphi=-\pi/2$: \textbf{droite} (sens horaire pour un observateur qui voit l'onde s'éloigner)
			      \item $\varphi=+\pi/2$: \textbf{gauche}
		      \end{itemize}
		\item \textbf{Elliptique}: cas général
		\item \textbf{Copolaire}: composante voulue; \textbf{croisée}: indésirable, à minimiser
		\item Antennes TX et RX doivent avoir la même polarisation (rectilignes: brins parallèles)
	\end{enumerate}
	\textbf{Désadaptation de polarisation} (perte en réception):

	\begin{tabular}{|l|l|}
		\hline
		\textbf{TX / RX}                       & \textbf{Perte}              \\
		\hline
		rectiligne / rectiligne, angle $\psi$ & $\cos^2\psi$ \\
		rectiligne / circulaire               & $\frac12$ ($-3$\,dB)        \\
		circulaire droite / gauche            & totale                      \\
		\hline
	\end{tabular}
	\begin{enumerate}[\roman*.]
		\item Brins croisés ($\psi=90^\circ$): rien n'est reçu
		\item Circulaire: insensible à la rotation de l'antenne (satellites)
	\end{enumerate}

	\colbreak
	\subsection{Antennes filaires}
	\textbf{Dipôle demi-onde}: 2 brins de $\lambda/4$ alimentés au centre.
	\begin{gather*}
		E_\theta=-j\,\frac{60I_0}{r}\,\frac{\cos(\frac\pi2\cos\theta)}{\sin\theta}e^{j(\omega t-kr)},\quad H_\varphi=\frac{E_\theta}{Z_0}\\
		D(\theta)=1{,}64\left(\frac{\cos(\frac\pi2\cos\theta)}{\sin\theta}\right)^2
	\end{gather*}
	\begin{enumerate}[\roman*.]
		\item $Z_{in}=73+j42\,\Omega$; polarisation rectiligne $\parallel$ au dipôle
		\item $D_{max}=1{,}64=2{,}15$\,dBi en $\theta=\pi/2$; nul selon l'axe ($\theta=0,\pi$)
		\item Indépendant de $\varphi$: \textbf{omnidirectionnel} dans le plan $\perp$ au dipôle; diagramme en «tore»; mêmes coupes $\varphi=0^\circ$ et $90^\circ$ (symétrie de révolution)
		\item Champs transverses, en $1/r$
	\end{enumerate}
	\textbf{Monopôle $\lambda/4$} au-dessus d'un plan de masse:
	\begin{enumerate}[\roman*.]
		\item Principe des images: équivaut à un dipôle demi-onde dans le demi-espace supérieur
		\item Directivité $+3$\,dB (toute la puissance dans un demi-espace)
		\item Plan fini: forme modifiée, symétrie de révolution conservée
		\item Usages ($\lambda/2$, $\lambda/4$): diffusion (TNT, radio), point-multipoint (Wi-Fi)
	\end{enumerate}
	\textbf{Dipôle replié}: 2 dipôles $\lambda/2$ reliés aux extrémités, proches: plus large bande.

	\textbf{Éléments parasites}: dipôles non alimentés, couplés: un peu plus long $\to$ \textbf{réflecteur}, plus court $\to$ \textbf{directeur}.

	\textbf{Yagi(-Uda)}: dipôle replié + 1 réflecteur + $N$ directeurs: directive, large bande (TV)
	\begin{enumerate}[\roman*.]
		\item $Z_{in}=75\,\Omega$; gain 4\,dB (1 directeur) à 20\,dB (20 directeurs)
		\item $\theta_{3dB}$: 60$^\circ$ (faible gain) à 25$^\circ$ (fort gain); AV/AR $>25$\,dB
	\end{enumerate}

	\colbreak
	\subsection{Ouvertures rayonnantes}
	Ouverture de surface $A$ éclairée par une onde (Huygens: sources secondaires):
	\begin{gather*}
		D_{max}=\frac{4\pi A}{\lambda^2}
	\end{gather*}
	\begin{enumerate}[\roman*.]
		\item \textbf{Cornet}: transition guide d'onde $\to$ espace libre; grand gain, faibles pertes, large bande
		\item \textbf{Réflecteur parabolique} (source au foyer): gain élevé, couverture modelable, multi-spots; encombrant
		      \begin{itemize}
			      \item Centrée: la source masque une partie du faisceau
			      \item Offset: meilleur rendement (TV satellite), réflecteur presque vertical
			      \item Cassegrain: second réflecteur hyperbolique, compact, plus directive
		      \end{itemize}
	\end{enumerate}
	Disque de diamètre $D_o$, efficacité $\eta$:
	\begin{gather*}
		D_{max}=\left(\frac{\pi D_o}{\lambda}\right)^2,\quad G=\eta D_{max},\quad\theta_{3dB}\approx70\frac{\lambda}{D_o}\ (^\circ)
	\end{gather*}
	\textbf{Exemple} 11\,GHz, $\eta=60\,\%$ ($-2{,}2$\,dB):

	\begin{tabular}{|l|l|l|l|}
		\hline
		$D_o$  & $D_{max}$  & $G$     & $\theta_{3dB}$ \\
		\hline
		1,8\,m & 46,3\,dBi & 44,1\,dB & 1$^\circ$      \\
		1,2\,m & 42,8\,dBi & 40,6\,dB & 1,6$^\circ$    \\
		\hline
	\end{tabular}

	\subsection{Réseaux et antennes imprimées}
	\begin{enumerate}[\roman*.]
		\item \textbf{Réseau}: plusieurs éléments, amplitudes/phases réglées (déphaseurs, diviseurs) pour augmenter le gain et \textbf{dépointer}: faisceau conformé, balayage électronique, adaptatif
		\item \textbf{Patch} (microruban): élément métallique + substrat + plan de masse; faible coût, encombrement, masse, conformable, intégrable; mais bande étroite, puissance limitée (claquage), gain $\sim6$\,dB; alimentation par microstrip ou coaxial
		\item \textbf{Monopôle imprimé}: sans plan de masse sous l'élément, se comporte comme une antenne filaire, omnidirectionnel
	\end{enumerate}

	\colbreak
	\subsection{Réflecteur plan (TD4, TD5)}
	Source à la hauteur $h$ d'un réflecteur de coefficient $\Gamma$: l'onde réfléchie a parcouru $2h$ de plus:
	\begin{gather*}
		\phi_R=2kh+\arg\Gamma,\quad\frac{p_{axe}}{p_{direct}}=\left|1+\Gamma e^{-2jkh}\right|^2
	\end{gather*}
	\textbf{Métallique} ($\Gamma=-1$): facteur $4\sin^2(kh)$
	\begin{enumerate}[\roman*.]
		\item Max pour $h=\lambda/4$: ondes en phase dans l'axe, antenne devenue directive (rayonnement arrière supprimé)
		\item Mieux que sans réflecteur si $\sin kh>\frac12$: $\frac\lambda{12}<h<\frac{5\lambda}{12}$, soit $\frac{c}{12h}<f<\frac{5c}{12h}$
		\item Inconvénient: $h\sim\lambda/4$, épais à basse fréquence; à $h\to0$ le réflecteur court-circuite l'antenne
	\end{enumerate}
	\textbf{Conducteur magnétique artificiel} (métamatériau, AMC: $\Gamma=+1$, impédance de surface infinie): facteur $4\cos^2(kh)$
	\begin{enumerate}[\roman*.]
		\item Réflexion en phase: fonctionne pour $h\to0$, pas de fréquence basse théorique; $f<\frac{c}{6h}$
		\item Avantage: antenne \textbf{mince} et unidirectionnelle (gain augmenté), impossible avec un métal
	\end{enumerate}

	\subsection{Méthode (lecture de fiche technique)}
	\begin{enumerate}[\roman*.]
		\item Courbe $|\Gamma|_{dB}(f)$: bande où la courbe est sous le seuil ($-10$\,dB ou ROS)
		\item Diagrammes: identifier le plan de coupe ($\varphi=0^\circ$: plan $xOz$; $\varphi=90^\circ$: $yOz$; $\theta=90^\circ$: $xOy$)
		\item Lire $G_{max}$, direction du max/min, $\theta_{3dB}$ (points à $G_{max}-3$\,dB), lobes secondaires, AV/AR
		\item Plaque métallique à $\lambda/4$: diagramme concentré d'un côté, $G_{max}\nearrow$
		\item Application: omni $\to$ diffusion/mobile; directive $\to$ point à point, satellite
	\end{enumerate}

	\colbreak
	\section{Bilan de liaison et propagation}
	\textbf{But:} prévoir la puissance reçue $P_R$ et la comparer à la \textbf{sensibilité} $P_{R,min}$ du récepteur (sinon signal inexploitable); en déduire portée, puissance d'émission, taille des antennes.

	\subsection{Équation de Friis}
	Liaison en espace libre, sans obstacle, antennes à la distance $d$, $e_{ray}=1$:
	\begin{gather*}
		p(d)=\frac{P_E}{4\pi d^2}G_E,\quad P_R=p\cdot S
	\end{gather*}
	\textbf{Surface équivalente} $S$ (surface de captation; $\leq$ ouverture physique):
	\begin{gather*}
		G_R=\frac{4\pi}{\lambda^2}S,\quad L_{FS}=\left(\frac{4\pi d}{\lambda}\right)^2\\
		\boxed{P_R=P_EG_EG_R\left(\frac{\lambda}{4\pi d}\right)^2}\\
		P_{R,dBm}=P_{E,dBm}+G_{E,dB}+G_{R,dB}-L_{FS,dB}\\
		L_{FS,dB}=20\log d_{km}+20\log f_{MHz}+32{,}45
	\end{gather*}
	\begin{enumerate}[\roman*.]
		\item $L_{FS}$ ($=$ PL, path loss): $+6$\,dB quand $d$ ou $f$ double
		\item Gains pris \textbf{dans la direction de la liaison} (lire le diagramme à l'angle concerné)
		\item Ajouter si besoin: désadaptation $-10\log(1-|\Gamma|^2)$, câbles, polarisation, pluie, gaz
	\end{enumerate}

	\subsection{PIRE}
	\textbf{Puissance isotrope rayonnée équivalente} (EIRP): l'émission; plafonnée par les normes.
	\begin{gather*}
		\mathrm{PIRE}=P_EG_E\ \text{(W)}\\ P_{R,dBW}=\mathrm{PIRE}_{dBW}+G_{R,dB}-L_{FS,dB}
	\end{gather*}
	\textbf{Exemple} (TD4) deux dipôles, $d=100$\,m, $P_E=0{,}5$\,W, 2,5\,GHz: dipôles \textbf{parallèles}, chacun face à l'autre dans son plan de max; $\lambda=12$\,cm, $L_{FS}=80{,}4$\,dB, $P_R=27+2{,}15+2{,}15-80{,}4\approx-49$\,dBm. Rotation de $30^\circ$: relire $G_R$ à $30^\circ$.

	\colbreak
	\subsection{Exemple satellite (TD4)}
	\textbf{Satellite géostationnaire}: $d=36\,000$\,km, 12\,GHz, $P_E=100$\,W (20\,dBW), $G_E=42$\,dB, $P_R=2\cdot10^{-11}$\,W ($-107$\,dBW) pour TEB $10^{-6}$:
	\begin{enumerate}[\roman*.]
		\item $\lambda=2{,}5$\,cm, $L_{FS}=205{,}2$\,dB
		\item $G_R=-107-20-42+205{,}2=36{,}2$\,dB
		\item $S=\frac{P_R\,4\pi d^2}{P_EG_E}=\frac{G_R\lambda^2}{4\pi}=0{,}21$\,m²
	\end{enumerate}

	\subsection{Facteurs de propagation}
	Réalité: absorption, réflexion, diffraction, diffusion, réfraction $\Rightarrow$ Friis optimiste.
	\begin{enumerate}[\roman*.]
		\item \textbf{Gaz} (dB/km): pics H$_2$O à 22,3; 183,3; 323\,GHz; O$_2$ à 60 et 118,74\,GHz; choix de $f$ selon portée longue ou communication discrète
		\item \textbf{Pluie}: atténuation $\nearrow$ avec $f$ et l'intensité; significative au-delà de 5\,GHz, prépondérante au-delà de 10\,GHz. Atténuation $=$ affaiblissement linéique $\times$ longueur du trajet sous la pluie; prévoir une marge
		\item \textbf{Sol}: conducteur sous $f_{transition}$ (mer $\sim$1\,GHz), diélectrique au-delà; image de l'antenne. \textbf{Diffraction} par une arête: liaison possible sans visibilité
	\end{enumerate}

	\subsection{Ellipsoïde de Fresnel}
	Visibilité directe radio $\neq$ optique: un obstacle proche de la ligne de visée diffracte, l'onde diffractée interfère avec l'onde directe. Règle: dégager le \textbf{premier ellipsoïde de Fresnel} (foyers: antennes; différence de trajet $\lambda/2$).
	\begin{gather*}
		\sqrt{d_1^2+r^2}+\sqrt{d_2^2+r^2}-(d_1+d_2)\approx\frac{r^2}2\Big(\frac1{d_1}+\frac1{d_2}\Big)=\frac\lambda2\\
		R_F=\sqrt{\lambda\frac{d_1d_2}{d_1+d_2}}\quad(d_1,d_2\gg r)
	\end{gather*}
	\begin{enumerate}[\roman*.]
		\item Max à mi-distance; $R_F\searrow$ quand $f\nearrow$ (vers l'optique)
		\item 1\,km, 2\,GHz, mi-chemin: 6\,m; TDF 8\,GHz, 57,6\,km: 23,2\,m
	\end{enumerate}

	\colbreak
	\textbf{Exemple} (TD4): antennes à $h_A=20$\,m, $d=10$\,km, 2,5\,GHz, sol plan:
	\begin{enumerate}[\roman*.]
		\item $d_1=d_2=5$\,km: $R_F=\sqrt{0{,}12\cdot2500}=17{,}3$\,m $<20$\,m: sol hors ellipsoïde
		\item Obstacle à $d_1=3$, $d_2=7$\,km: $R_F=15{,}9$\,m $\Rightarrow h_{max}=20-15{,}9=4{,}1$\,m
		\item $P_E=5$\,W (37\,dBm), $G=20$\,dB chacune: $L_{FS}=120{,}4$\,dB, $P_R=-43{,}4$\,dBm $=46$\,nW
	\end{enumerate}

	\subsection{Modèles empiriques}
	Milieu urbain/rural sans visibilité directe, multi-trajets: modèles issus de mesures (Hata, Okumura-Hata, 3GPP\dots).

	\textbf{Okumura-Hata rural} ($h_m\leq1{,}5$\,m; $f$ en MHz, $h_{bs}$ en m, $d$ en km):
	\begin{gather*}
		PL=69{,}55+26{,}16\log f-13{,}82\log h_{bs}\\
		+(44{,}9-6{,}55\log h_{bs})\log d\\
		-4{,}78(\log f)^2+18{,}33\log f-40{,}94
	\end{gather*}
	\textbf{Exemple} GSM: 948\,MHz, $d=20$\,km, $h_{bs}=130$\,m, $P_E=8$\,W (39\,dBm), $G_E=14$, $G_R=0$\,dB:
	\begin{enumerate}[\roman*.]
		\item Hata: $PL=129{,}9$\,dB, $P_R=-76{,}8$\,dBm
		\item Friis: $L_{FS}=118{,}0$\,dB, $P_R=-65{,}0$\,dBm (optimiste de 12\,dB)
		\item Distance Friis équivalente: $d'=d\cdot10^{(129{,}9-118{,}0)/20}=78{,}5$\,km
	\end{enumerate}

	\subsection{Méthode}
	\begin{enumerate}[\roman*.]
		\item Tout en dB: $P$ (dBm), $G$ (dB), $L$ (dB); $\lambda$ d'abord
		\item $P_R=P_E+G_E+G_R-L$; comparer à la sensibilité, la marge restante absorbe pluie et gaz
		\item Inconnue $G_R$ ou $S$: isoler dans Friis puis $S=G\lambda^2/4\pi$
		\item Dégagement: $R_F$ au point de l'obstacle, comparer à la hauteur libre
	\end{enumerate}

	%%%%%%%%%%%%%%%%%%%%%%%%%%%%%%%%%%%%%%%%%%%%%%%%%%%%%%
	%                       End                          %
	%%%%%%%%%%%%%%%%%%%%%%%%%%%%%%%%%%%%%%%%%%%%%%%%%%%%%%

\end{multicols*}
\end{document}
